\documentclass[10pt, letterpaper]{article}

% Packages:
\usepackage[
    ignoreheadfoot,
    top=0.4 cm,
    bottom=0.4 cm,
    left=0.8 cm,
    right=0.8 cm,
    footskip=0.2 cm,
]{geometry}
\usepackage{titlesec}
\usepackage{tabularx}
\usepackage{array}
\usepackage[dvipsnames]{xcolor}
\definecolor{primaryColor}{RGB}{0, 0, 0}
\usepackage{enumitem}
\usepackage{amsmath}
\usepackage[
    pdftitle={Tanmay Bhuskute's CV},
    pdfauthor={Tanmay Bhuskute},
    pdfcreator={LaTeX with RenderCV},
    colorlinks=true,
    urlcolor=primaryColor
]{hyperref}
\usepackage[pscoord]{eso-pic}
\usepackage{calc}
\usepackage{bookmark}
\usepackage{lastpage}
\usepackage{changepage}
\usepackage{paracol}
\usepackage{ifthen}
\usepackage{needspace}
\usepackage{iftex}

\ifPDFTeX
    \input{glyphtounicode}
    \pdfgentounicode=1
    \usepackage[T1]{fontenc}
    \usepackage[utf8]{inputenc}
    \usepackage{lmodern}
\fi

\usepackage{XCharter}

% Some settings:
\raggedright
\AtBeginEnvironment{adjustwidth}{\partopsep0pt}
\pagestyle{empty}
\setcounter{secnumdepth}{0}
\setlength{\parindent}{0pt}
\setlength{\topskip}{0pt}
\setlength{\columnsep}{0.15cm}
\pagenumbering{gobble}

% ── Section title ──
\titleformat{\section}{\needspace{2\baselineskip}\bfseries\large}{}{0pt}{}[\vspace{1pt}\titlerule]
\titlespacing{\section}{-1pt}{0.15 cm}{0.08 cm}

\renewcommand\labelitemi{$\vcenter{\hbox{\small$\bullet$}}$}

% ── Bullet lists ──
\newenvironment{highlights}{
    \begin{itemize}[
        topsep=0.10 cm,
        parsep=0.06 cm,
        partopsep=0pt,
        itemsep=0.06 cm,
        leftmargin=12pt
    ]
}{
    \end{itemize}
}

\newenvironment{highlightsforbulletentries}{
    \begin{itemize}[
        topsep=0.10 cm,
        parsep=0.06 cm,
        partopsep=0pt,
        itemsep=0.04 cm,
        leftmargin=12pt
    ]
}{
    \end{itemize}
}

\newenvironment{onecolentry}{
    \begin{adjustwidth}{0 cm + 0.00001 cm}{0 cm + 0.00001 cm}
}{
    \end{adjustwidth}
}

\newenvironment{twocolentry}[2][]{
    \onecolentry
    \def\secondColumn{#2}
    \setcolumnwidth{\fill, 4.5 cm}
    \begin{paracol}{2}
}{
    \switchcolumn \raggedleft \secondColumn
    \end{paracol}
    \endonecolentry
}

\newenvironment{threecolentry}[3][]{
    \onecolentry
    \def\thirdColumn{#3}
    \setcolumnwidth{, \fill, 4.5 cm}
    \begin{paracol}{3}
    {\raggedright #2} \switchcolumn
}{
    \switchcolumn \raggedleft \thirdColumn
    \end{paracol}
    \endonecolentry
}

\newenvironment{header}{
    \setlength{\topsep}{0pt}\par\kern\topsep\centering\linespread{0.6}
}{
    \par\kern\topsep
}

\let\hrefWithoutArrow\href

\begin{document}
    \newcommand{\AND}{\unskip
        \cleaders\copy\ANDbox\hskip\wd\ANDbox
        \ignorespaces
    }
    \newsavebox\ANDbox
    \sbox\ANDbox{$|$}

    \begin{header}
        \fontsize{16 pt}{24 pt}\selectfont Tanmay Bhuskute

        \vspace{3 pt}

        \normalsize
        \mbox{ +1 (732) 858-3027 }%
        \kern 1.0 pt%
        \AND%
        \kern 1.0 pt%
        \mbox{ tanmay.v.bhuskute@gmail.com }%
        \kern 1.0 pt%
        \AND%
        \kern 1.0 pt%
        \mbox{ github.com/Tanmayb05 }%
        \kern 1.0 pt%
        \AND%
        \kern 1.0 pt%
        \mbox{ linkedin.com/in/tanmay-bhuskute }%
        \kern 1.0 pt%
        \AND%
        \kern 1.0 pt%
        \mbox{ Cincinnati, OH }%
    \end{header}

    \vspace{2 pt}

%═══════════════════════════════════════
\section{SUMMARY}
%═══════════════════════════════════════
\begin{onecolentry}
Software Engineer and M.S. Computer Science candidate (Dec 2026) with production experience designing, building, testing, and operating backend services and cloud infrastructure using Java, Python, Go, and TypeScript. Strong in REST API design, relational database schema and migrations, microservices, CI/CD testing, and observability.
\end{onecolentry}

%═══════════════════════════════════════
\section{EDUCATION}
%═══════════════════════════════════════
\begin{twocolentry}{
    Aug 2024 - Dec 2026
}
    \textbf{Master of Science in Computer Science} \\
    University of Cincinnati
\end{twocolentry}
\begin{onecolentry}
Relevant Coursework: Advanced Algorithms, Database Theory, Intelligent Data Analysis.
\end{onecolentry}
\vspace{0.10 cm}

\begin{twocolentry}{
    Aug 2018 - Jul 2022
}
    \textbf{Bachelor of Engineering in Computer Engineering} \\
    AISSMS College of Engineering
\end{twocolentry}
\begin{onecolentry}
Relevant Coursework: Object-Oriented Programming, Data Analytics, Machine Learning.
\end{onecolentry}
%═══════════════════════════════════════
\section{TECHNICAL SKILLS}
%═══════════════════════════════════════
\begin{onecolentry}
\textbf{Programming:} Java, Python, Go, TypeScript, JavaScript, C/C++, SQL, C\#.
\end{onecolentry}
\vspace{0.06 cm}
\begin{onecolentry}
\textbf{Backend \& Data:} Node.js, FastAPI, Flask, Express, REST APIs, GraphQL, Microservices, Message Queues (SQS); PostgreSQL, MySQL, Redis, schema design, indexing, query tuning, migrations.
\end{onecolentry}
\vspace{0.06 cm}
\begin{onecolentry}
\textbf{Cloud \& DevOps:} AWS (EC2, S3, Lambda, SQS, RDS, CloudWatch, IAM), Docker, Kubernetes, Terraform, GitLab CI/CD, Jenkins, ArgoCD, Helm, Linux, SonarQube, Datadog.
\end{onecolentry}
\vspace{0.06 cm}
\begin{onecolentry}
\textbf{Practices:} Object-Oriented Design, SOLID, Design Patterns, TDD, Unit/Integration/E2E Testing, Code Review, Agile/Scrum.
\end{onecolentry}

%═══════════════════════════════════════
\section{WORK EXPERIENCE}
%═══════════════════════════════════════
\begin{twocolentry}{Jul 2022 -- Jun 2024}
\textbf{Software Development Engineer} \,|\, Siemens Industry Software
\end{twocolentry}
\begin{highlights}
  \item Designed and maintained production backend automation for \textbf{Teamcenter X} in Python (Boto3) with object-oriented service modules and REST API integrations, replacing legacy Ruby scripts and cutting deployment time by \textbf{40\%} across AWS environments.
  \item Migrated \textbf{40+ production environments} from process-based deployments to containerized \textbf{Kubernetes} microservices using ArgoCD, Helm, GitLab CI/CD, and HashiCorp Vault, an incremental legacy-modernization effort with backward-compatible cutover.
  \item Modeled environment config, provisioning state, and job status in \textbf{DynamoDB} and relational stores with indexed lookups, supporting idempotent retries and partial-failure recovery in CLI and Ansible workflows; added automated rollback cutting upgrade downtime \textbf{50\%}.
  \item Executed \textbf{44 production environment upgrades} in 4 months, triaging live incidents across Linux, EC2, and databases using CloudWatch metrics, tracing, and Datadog structured-log correlation; as Quality Coach and \textbf{SonarQube} Specialist raised automated test coverage from \textbf{35\% to 78\%} via CI/CD quality gates and mentored \textbf{5 engineers}.
\end{highlights}

\vspace{0.10 cm}

\begin{twocolentry}{Feb 2026 -- Present}
\textbf{Graduate Researcher} \,|\, University of Cincinnati
\end{twocolentry}
\begin{highlights}
  \item Designed two-stage memory-augmented LLM pipeline for Android GUI bug reproduction as a config-driven, provider-agnostic service with a versioned run-artifact schema; evaluated against \textbf{ViBR} baseline across 53 runs on 20+ apps, showing \textbf{25.7x latency reduction} while preserving task-completion fidelity.
\end{highlights}

\vspace{0.10 cm}

\begin{twocolentry}{Mar 2022 -- Jun 2022}
\textbf{Software Developer Intern} \,|\, Siemens Industry Software
\end{twocolentry}
\begin{highlights}
  \item Automated AWS provisioning with \textbf{Terraform} infrastructure-as-code (\textbf{30\%} faster deploys) and built a \textbf{Go + Terratest} framework catching \textbf{95\%} of misconfigurations pre-production in CI.
\end{highlights}

\vspace{0.10 cm}

\begin{twocolentry}{Sep 2021 -- Dec 2021}
\textbf{Software Developer Intern} \,|\, Government of Maharashtra, Water Resource Department
\end{twocolentry}
\begin{highlights}
  \item Built and deployed a full-stack \textbf{React and MySQL} workflow application that digitized departmental record management and improved operational efficiency by \textbf{40\%}.
\end{highlights}

%═══════════════════════════════════════
\section{PROJECTS AND RESEARCH}
%═══════════════════════════════════════
\begin{twocolentry}{May 2026}
\textbf{PolicyPilot -- Policy Intelligence Portal} [React, TS, Node, Express, Gemini, MCP]
\end{twocolentry}
\begin{highlights}
  \item Built full-stack TypeScript app (React SPA + \textbf{Express 5 REST API} + tool server) on one Node.js process with \textbf{17 versioned tool APIs} and Zod schema validation at every data boundary; tested with \textbf{Vitest} suites in CI.
\end{highlights}

\vspace{0.10 cm}

\begin{twocolentry}{Jul 2026}
\textbf{NYC Mobility Forecasting} [Python, FastAPI, PostgreSQL, Redis, Docker]
\end{twocolentry}
\begin{highlights}
  \item Built \textbf{FastAPI REST service} serving real-time demand predictions across \textbf{263 zones} from a \textbf{2.3M-row} dataset, backed by PostgreSQL with a \textbf{Redis} cache and containerized with Docker.
\end{highlights}

%═══════════════════════════════════════
\section{EXTRACURRICULAR ACTIVITIES}
%═══════════════════════════════════════
\begin{onecolentry}
\begin{highlights}
  \item \textbf{Finalist, RevolutionUC 2025 Hackathon} (300 participants, 31 schools) \,|\, \textbf{SAFe Practitioner 5.1} (Mar 2023).
\end{highlights}
\end{onecolentry}

\end{document}
